% Part: incompleteness
% Chapter: theories-computability
% Section: introduction

\documentclass[../../../include/open-logic-section]{subfiles}

\begin{document}

\olfileid{inc}{tcp}{int}

\olsection{引言}

\begin{editorial}
本节需重写。
\end{editorial}

我们已有以下内容：
\begin{enumerate}
\item 函数在 $\Th{Q}$ 中可表示的定义（\olref[req][int]{defn:representable-fn}）
\item 关系在 $\Th{Q}$ 中可表示的定义（\olref[req][rel]{defn:representing-relations}）
\item 断言 $\Th{Q}$ 的可表示函数恰好就是可计算函数的定理（\olref[req][int]{thm:representable-iff-comp}）
\item 断言 $\Th{Q}$ 的可表示关系恰好就是可计算关系的定理（\olref[req][rel]{thm:representing-rels}）
\end{enumerate}

{\em 理论}是演绎封闭的语句集合，即：只要 $T$ 证明了 $!A$，$!A$ 就属于 $T$。最好将理论理解为一组语句，连同由这些语句所蕴涵的一切。从现在起，我们将用 $\Th{Q}$ 指称由\olref[req][int]{sec}中八条公理可推导出的全体语句所构成的{\em 理论}。请注意，我们可以将 $\Th{Q}$ 的公式编码为数字；若 $!A$ 是这样一个公式，令 $\Gn{!A}$ 表示编码 $!A$ 的数字。借助于此编码，我们现在就可以探讨各种公式集合是否可计算。

\end{document}